% =====================================================================
%  sections/sec8.tex —— 8 结论 + 致谢（原文 pp.48–50）
% =====================================================================
\section{结论}
\label{sec:concl}

本文用六维构造造出了一大类 AD 型理论。我们求出了 Seiberg--Witten 曲线与谱的标度维数，还借助三维镜像对称算出了中心荷 $a$ 与 $c$。我们也识别出了 $\SU(N)$ QCD 的 AD 点。

最初，有拉氏量描述的 $\sN{=}2$ SCFT 由 ADE Dynkin 图分类 \cite{49}。Gaiotto 用黎曼面上的正则穿刺点找到了一大类新的 SCFT \cite{2}，而这些理论一般没有拉氏量描述。在球面的情形，它们就是带强耦合物质系统的 $A_n$ 型广义箭图规范理论。这些结果极大地增加了 $\sN{=}2$ SCFT 的数目。上述理论都只有整数标度维数与无量纲的耦合常数。通过允许分数标度维数，我们说明了 $\sN{=}2$ SCFT 的版图被极大扩展，而且实际上把此前构造的那些理论（至少是定义在球面上的那些）作为一个很小的角落包含进来。考虑到 SCFT 的重要性，如此大量的例子应当对我们认识量子场论的动力学非常有帮助。

AD 理论虽然在 1995 年就被发现了，但人们对它们的性质了解甚少，相关研究也不多。有了六维构造，我们相信现在可以更细致地理解这些理论。大体上说，凡是针对 class S 理论做过的事情，稍加努力都可以推广到这些 AD 理论上。下面列出一些有意思的问题：

\begin{enumerate}[label=\textup{\alph*.}]
  \item 本文只考虑了六维 $A_n$ 理论，推广到六维 $D_n$ 与 $E_n$ 理论应当是直接的。文献 \cite{28} 研究了一般约化群的非正则奇点与广义单值性，可以利用这些结果构造新类型的 Argyres--Douglas 理论。
  \item 四维 $\sN{=}2$ 理论引力对偶的一般拟设已在 \cite{50} 中写出，而引力解归结为求解三维 Toda 方程。正则奇点的情形已在 \cite{51} 中研究过。通过研究 Toda 方程的非正则解，应当能够找到本文所构造 AD 理论的引力对偶。由于定义在球面上的 class S 理论也有非正则奇点的表示，利用非正则奇点或许更容易得到引力对偶。高秩 $(A_0,A_1,A_2)$ AD 理论的引力对偶已在 \cite{52} 中用 F 理论找到。既然我们给出了显式的六维构造，上述结果也许有助于求出一般的解。
  \item 只由正则穿刺点定义的场论，其丛坐标已在 \cite{53} 中构造出来。那套构造同样适用于本文构造的 AD 理论。思路是这样的：非正则奇点造出一个带多个标记点的边界，因此几何其实是带标记点的圆盘，\cite{53} 的构造很容易搬过来。问题就归结为：对某个特定的非正则奇点，要确定边界上有多少个标记点、以及该放什么杨图 \cite{30}。一旦求出丛坐标，就可以用它们研究线算子的分类、BPS 谱与墙穿越、面算子等等。
  \item $A_1$ AD 理论的 AGT 猜想 \cite{54} 已在 \cite{55,56} 中研究过，其中牵涉 Liouville 理论的非正则共形块。把这个猜想推广到本文构造的更高秩 AD 理论应当是可以做到的。由于这些理论没有拉氏量，必须另想办法求 Nekrasov 配分函数；在存在非正则奇点时，只要我们对平坦联络的模空间有好的坐标，\cite{57} 所述的方法应当有用。由 $A_1$ 理论构造的渐近自由理论的 AGT 猜想也已有研究 \cite{58}，把那一构造推广到本文构造的所有渐近自由理论同样有意思。
  \item Seiberg--Witten 解定义一个可积系统，而四维 $\sN{=}2$ 理论的配分函数与可积系统的 Yang--Yang 函数之间的 Nekrasov--Shatashvili 关系，在这里也可以推广。看起来，只要丛坐标已知，\cite{59} 构造的 Darboux 坐标就可以推广到我们的情形。事实上，在正则穿刺点的情形，他们用的是裤分解中的那些圈，而这些圈实际上标定了 Teichmüller 空间的叶层化空间。把叶层化推广到圆盘的做法见 \cite{60}；只要知道丛坐标，就应当能够写出 Darboux 坐标对应表达式的描述。
  \item 数学上，非正则奇点与可积系统近年来已被大量研究 \cite{61}，我们预期这些结果对计算 AD 理论的共形块有用。我们已经说明，同一个理论既可以用正则奇点、也可以用非正则奇点来实现；这一事实对研究几何 Langlands 纲领中的驯顺分歧与野分歧是否有用？此外，两种实现对应的等单值形变方程已被证明是等价的 \cite{62}，看看这对相应的场论意味着什么是很有意思的。
  \item 一旦知道这些更高秩 AD 理论的 BPS 谱，就可以沿 \cite{63} 中发展出的路子研究更高秩的三维四面体理论与相应的 $\sN{=}2$ 理论。例如，更高秩五边形理论的 BPS 谱似乎比较容易算（$A_2$ 实现的 $E_7$ 理论边界上有五个标记点）。
  \item 我们已经算出了一些 AD 理论的中心荷 $a$ 与 $c$。在这些理论之间找出重整化群流并不困难，显式检验 $a$ 定理应当很有意思。四维 $a$ 定理的一般证明在 \cite{64} 中讨论过；若能构造许多显式例子，对理解那个证明应当有帮助。
  \item 对每个 AD 理论，它的丛坐标对应于圆盘上的一个平面二部图。这类平面二部图在研究 $\sN{=}4$ 超杨–米尔斯理论的散射振幅时被大量使用 \cite{65}。同一个数学结构出现在两个看来完全不同的物理系统中，弄清这里是否有更深层的联系很有意思。AD 理论的物理可观测量与 $\sN{=}4$ 超杨–米尔斯的散射振幅会有关系吗？此外，平面二部图已在 \cite{66} 中被完全分类，能否用这个结果给 AD 理论一个另外的分类？
  \item 还有一些别的物理量是人们想算的。例如：算出 type IV AD 理论的中心荷 $k$ 会很有意思，它衡量非阿贝尔味对称性被规范之后对 $\beta$ 函数的贡献。算这些理论的超共形指标也有意思，正则穿刺点情形的结果见 \cite{67}。我们还想更细致地研究 AD 理论的希格斯分支，它似乎可以由三维镜像理论的库仑分支来捕捉，如此等等。
\end{enumerate}

\typo{第 (d) 条里的「AGT 猜想」在原文 p.49 有一处印作 "ACT conjecture"（同一条前面写的是 AGT），应为 AGT。第 (j) 条原文用 $k$ 表示这个中心荷，与全文用 $k$ 表示 block 大小、规范群秩的用法冲突；译文照录。}

\section*{致谢}
\addcontentsline{toc}{section}{致谢}

感谢 Nima Arkani-Hamed、Davide Gaiotto、Yu-tin Huang、Hai Lin、Juan Maldacena、Andrew Neitzke、Edward Witten、Yuji Tachikawa、Peng Zhao 的有益讨论。DX 的研究受 Zurich Financial Services 的会员资格资助，并感谢美国能源部 grant DE-FG02-90ER40542 的支持。
